%%%%%%%%%%%%%%%%%%%%%%%%%%%%%%%%%%%%%%%%%%%%%%%%%%%%%%%%%%%%%%%%%%%%%%%%%%%%%%%
%% TKS v7.4 SEMANTICS TRACK --- Domain-Theoretic Semantics
%% Part of TKS v7.4 - December 2025
%% ADDITIVE to v7.3 canon
%%%%%%%%%%%%%%%%%%%%%%%%%%%%%%%%%%%%%%%%%%%%%%%%%%%%%%%%%%%%%%%%%%%%%%%%%%%%%%%

%% ============================================================================
%% CHAPTER 1: SCOTT DOMAINS FOR NOETIC SPACES
%% ============================================================================

\chapter{Scott Domains for Noetic Spaces}
\label{ch:scott-domains-noetic}

\begin{tcolorbox}[colback=blue!5!white,colframe=blue!75!black,title=\vnewfour]
This chapter develops Scott domains as the semantic foundation for noetic
computation, establishing that the space of Ideas admits a natural domain structure.
\end{tcolorbox}

\section{Preliminaries: Domain Theory Review}

\begin{definition}[Directed Set]
\label{def:directed-set}
A subset $D \subseteq P$ of a poset $(P, \sqleq)$ is \textbf{directed} if:
\begin{enumerate}[label=(\roman*)}
  \item $D \neq \varnothing$ (nonemptiness), and
  \item For all $x, y \in D$, there exists $z \in D$ with $x \sqleq z$ and $y \sqleq z$.
\end{enumerate}
\end{definition}

\begin{definition}[Directed-Complete Partial Order (DCPO)]
\label{def:dcpo}
A poset $(D, \sqleq)$ is a \textbf{dcpo} if every directed subset has a supremum.
\end{definition}

\begin{definition}[Scott Continuity]
\label{def:scott-continuous}
A function $f : D \to E$ between dcpos is \textbf{Scott-continuous} if:
\begin{enumerate}[label=(\roman*)]
  \item $f$ is monotone: $x \sqleq y \Rightarrow f(x) \sqleq f(y)$, and
  \item $f$ preserves directed suprema: $f\left(\bigsqcup S\right) = \bigsqcup \{f(s) \mid s \in S\}$
\end{enumerate}
\end{definition}

\section{Noetic Domains}

\begin{definition}[Noetic Domain]
\label{def:noetic-domain}
The \textbf{noetic domain} $\NoeticDomain$ is a Scott domain constructed as:
\begin{enumerate}[label=(\roman*)]
  \item \textbf{Base elements}: The Ideas $\Ideas = \{X_1, \ldots, X_{40}\}$ form maximal elements.
  \item \textbf{Approximation structure}: Define $\Ideas_\bot = \Ideas \cup \{\bot_\nu\}$.
  \item \textbf{Order}: The flat domain order with $\bot_\nu$ below all Ideas.
\end{enumerate}
\end{definition}

\begin{proposition}[Flat Domain Properties]
\label{prop:flat-domain-props}
The noetic domain $\NoeticDomain = \Ideas_\bot$ satisfies:
\begin{enumerate}[label=(\roman*)]
  \item $\NoeticDomain$ is a pointed dcpo with bottom $\bot_\nu$.
  \item Every element except $\bot_\nu$ is maximal.
  \item Every directed set is either $\{\bot_\nu\}$ or eventually constant.
  \item $\NoeticDomain$ is algebraic with $K(\NoeticDomain) = \NoeticDomain$.
\end{enumerate}
\end{proposition}

\begin{definition}[Extended Noetic Domain]
\label{def:extended-noetic-domain}
The \textbf{extended noetic domain} $\NoeticDomain^\infty$ incorporates transfinite
fractal structure:
\[
\NoeticDomain^\infty = \NoeticDomain \times \OrdFrac_\bot
\]
with the product order.
\end{definition}

\begin{theorem}[Extended Noetic Domain is a DCPO]
\label{thm:extended-noetic-dcpo}
The extended noetic domain $\NoeticDomain^\infty$ is a pointed dcpo.
\end{theorem}

\section{Algebraicity and Continuous Domains}

\begin{definition}[Compact Element]
\label{def:compact-element}
An element $c \in D$ of a dcpo is \textbf{compact} if for every directed set $S$
with $c \sqleq \bigsqcup S$, there exists $s \in S$ with $c \sqleq s$.
\end{definition}

\begin{definition}[Scott Domain]
\label{def:scott-domain}
A \textbf{Scott domain} is an algebraic dcpo that is bounded complete and
$\omega$-algebraic.
\end{definition}

\begin{theorem}[Noetic Domain is a Scott Domain]
\label{thm:noetic-scott-domain}
The noetic domain $\NoeticDomain$ is a Scott domain.
\end{theorem}

\section{Function Spaces and Exponentials}

\begin{theorem}[Function Space is a DCPO]
\label{thm:function-space-dcpo}
If $D$ and $E$ are dcpos, then $[D \to E]$ is a dcpo with pointwise order.
\end{theorem}

\begin{corollary}[Noetic Function Domain]
\label{cor:noetic-function-domain}
The space $[\NoeticDomain \to \NoeticDomain]$ is a pointed dcpo containing
all noetic operators $\noe{k}$.
\end{corollary}

%% ============================================================================
%% CHAPTER 2: CONTINUOUS LATTICES AND WAY-BELOW
%% ============================================================================

\chapter{Continuous Lattices and Way-Below}
\label{ch:continuous-lattices}

\begin{tcolorbox}[colback=blue!5!white,colframe=blue!75!black,title=\vnewfour]
This chapter introduces the way-below relation $\WayBelow$, capturing finite
approximation and leading to continuous lattice structure.
\end{tcolorbox}

\section{The Way-Below Relation}

\begin{definition}[Way-Below Relation]
\label{def:way-below}
For elements $x, y$ of a dcpo, we say $x$ is \textbf{way below} $y$, written
$x \WayBelow y$, if for every directed set $S$ with $y \sqleq \bigsqcup S$,
there exists $s \in S$ with $x \sqleq s$.
\end{definition}

\begin{lemma}[Way-Below Properties]
\label{lem:way-below-props}
The way-below relation satisfies:
\begin{enumerate}[label=(\roman*)]
  \item $x \WayBelow y \Rightarrow x \sqleq y$
  \item $x' \sqleq x \WayBelow y \sqleq y' \Rightarrow x' \WayBelow y'$
  \item $\bot \WayBelow y$ for all $y$ (if $D$ is pointed)
  \item Transitivity with order
\end{enumerate}
\end{lemma}

\begin{proposition}[Way-Below in Flat Domains]
\label{prop:way-below-flat}
In $\NoeticDomain$:
\[
x \WayBelow y \quad \Leftrightarrow \quad x = \bot_\nu \text{ or } x = y
\]
\end{proposition}

\section{Continuous Lattices}

\begin{definition}[Continuous DCPO]
\label{def:continuous-dcpo}
A dcpo $D$ is \textbf{continuous} if for every $x \in D$:
\[
x = \bigsqcup \{y \in D \mid y \WayBelow x\}
\]
and this set is directed.
\end{definition}

\begin{definition}[Continuous Lattice]
\label{def:continuous-lattice}
A \textbf{continuous lattice} is a complete lattice that is continuous as a dcpo.
\end{definition}

\begin{theorem}[Noetic Domain is Continuous]
\label{thm:noetic-continuous}
The noetic domain $\NoeticDomain$ is a continuous dcpo.
\end{theorem}

\section{Interpolation Properties}

\begin{theorem}[Continuous DCPOs Have Interpolation]
\label{thm:continuous-interpolation}
Every continuous dcpo has the interpolation property: if $x \WayBelow z$,
there exists $y$ with $x \WayBelow y \WayBelow z$.
\end{theorem}

%% ============================================================================
%% CHAPTER 3: FIXED-POINT SEMANTICS
%% ============================================================================

\chapter{Fixed-Point Semantics}
\label{ch:fixed-point-semantics}

\begin{tcolorbox}[colback=blue!5!white,colframe=blue!75!black,title=\vnewfour]
This chapter develops fixed-point semantics for recursive noetic operations
via the Kleene fixed-point theorem.
\end{tcolorbox}

\section{Kleene Fixed-Point Theorem}

\begin{theorem}[Kleene Fixed-Point Theorem]
\label{thm:kleene-fixpoint}
Let $(D, \sqleq, \bot)$ be a pointed dcpo and $f : D \to D$ Scott-continuous.
Then $f$ has a least fixed point:
\[
\KleeneFix(f) = \bigsqcup_{n < \omega} f^n(\bot)
\]
\end{theorem}

\begin{proof}
\textbf{Step 1}: The sequence $\{f^n(\bot)\}$ is ascending by induction.

\textbf{Step 2}: The sequence is directed (ascending implies directed).

\textbf{Step 3}: The supremum $x^* = \bigsqcup_n f^n(\bot)$ exists by dcpo property.

\textbf{Step 4}: $x^*$ is a fixed point by continuity:
\[
f(x^*) = f\left(\bigsqcup_n f^n(\bot)\right) = \bigsqcup_n f^{n+1}(\bot) = x^*
\]

\textbf{Step 5}: $x^*$ is least: any fixed point $y$ satisfies $f^n(\bot) \sqleq y$
for all $n$, hence $x^* \sqleq y$.
\end{proof}

\begin{theorem}[Continuity of $\KleeneFix$]
\label{thm:lfp-continuous}
The operator $\KleeneFix : [D \to D] \to D$ is Scott-continuous.
\end{theorem}

\section{Noetic Fixed Points}

\begin{definition}[Noetic Fixed Point]
\label{def:noetic-fixed-point}
For a Scott-continuous noetic transformation $F : \NoeticDomain \to \NoeticDomain$:
\[
\fix_\nu(F) = \KleeneFix(F) = \bigsqcup_{n < \omega} F^n(\bot_\nu)
\]
\end{definition}

\begin{proposition}[Fixed Points in Flat Domains]
\label{prop:flat-fixed-points}
In $\NoeticDomain$, the Kleene chain stabilizes in at most two steps:
\[
\KleeneFix(F) =
\begin{cases}
\bot_\nu & \text{if } F(\bot_\nu) = \bot_\nu \\
F(\bot_\nu) & \text{otherwise}
\end{cases}
\]
\end{proposition}

\section{Tarski-Knaster Fixed Points}

\begin{theorem}[Tarski-Knaster Fixed Point Theorem]
\label{thm:tarski-knaster}
Let $L$ be a complete lattice and $f : L \to L$ monotone. Then:
\begin{enumerate}[label=(\roman*)]
  \item The set $\mathrm{Fix}(f)$ of fixed points forms a complete lattice.
  \item The least fixed point is $\mu f = \bigwedge\{x \mid f(x) \sqleq x\}$.
  \item The greatest fixed point is $\nu f = \bigvee\{x \mid x \sqleq f(x)\}$.
\end{enumerate}
\end{theorem}

\section{Recursive Noetic Definitions}

\begin{definition}[Recursive Noetic Specification]
\label{def:recursive-spec}
A \textbf{recursive noetic specification} is an equation $X = F(X)$ where
$F : \NoeticDomain \to \NoeticDomain$ is Scott-continuous.
The denotation is $\sem{X = F(X)} = \KleeneFix(F)$.
\end{definition}

\begin{example}[Infinite Fractal Application]
\label{ex:infinite-fractal}
The specification $X = \noe{k}(X)$ asks for an Idea that is a fixed point of $\noe{k}$.
\end{example}

%% ============================================================================
%% CHAPTER 4: DOMAIN EQUATIONS FOR NOETIC STRUCTURES
%% ============================================================================

\chapter{Domain Equations for Noetic Structures}
\label{ch:domain-equations}

\begin{tcolorbox}[colback=blue!5!white,colframe=blue!75!black,title=\vnewfour]
This chapter addresses domain equations $D \DomainEq F(D)$ for modeling
recursive noetic types.
\end{tcolorbox}

\section{Recursive Domain Equations}

\begin{definition}[Domain Equation]
\label{def:domain-equation}
A \textbf{domain equation} is an isomorphism specification:
\[
D \DomainEq F(D)
\]
where $F : \ScottDom \to \ScottDom$ is an endofunctor.
\end{definition}

\begin{example}[Canonical Domain Equations]
\label{ex:canonical-domain-eq}
\begin{enumerate}[label=(\alph*)]
  \item \textbf{Lazy lists}: $L \cong 1 + A \times L$
  \item \textbf{Binary trees}: $T \cong 1 + T \times T$
  \item \textbf{Streams}: $S \cong A \times S$
  \item \textbf{Functions}: $D \cong [D \to D]$
\end{enumerate}
\end{example}

\section{Solving Domain Equations}

\begin{theorem}[Existence of Solutions]
\label{thm:domain-eq-existence}
Let $F : \ScottDom \to \ScottDom$ be locally continuous and preserve
embedding-projection pairs. Then $D \cong F(D)$ has a solution, unique up to isomorphism.
\end{theorem}

\begin{construction}[Bilimit Construction]
\label{constr:bilimit}
The solution is constructed via inverse limits:
\[
D_\infty = \varprojlim_n D_n
\]
where $D_0 = \{\bot\}$ and $D_{n+1} = F(D_n)$.
\end{construction}

\section{Initial Algebra Semantics}

\begin{definition}[$F$-Algebra]
\label{def:f-algebra}
For an endofunctor $F$, an \textbf{$F$-algebra} is a pair $(A, \alpha)$ with
$\alpha : F(A) \to A$.
\end{definition}

\begin{theorem}[Lambek's Lemma]
\label{thm:lambek}
If $(I, \iota)$ is an initial $F$-algebra, then $\iota : F(I) \to I$ is an isomorphism.
\end{theorem}

\section{Final Coalgebra Duality}

\begin{definition}[$F$-Coalgebra]
\label{def:f-coalgebra}
An \textbf{$F$-coalgebra} is a pair $(C, \gamma)$ with $\gamma : C \to F(C)$.
\end{definition}

\begin{theorem}[Final Coalgebra is Fixed Point]
\label{thm:final-coalgebra-fixpoint}
If $(Z, \zeta)$ is a final $F$-coalgebra, then $\zeta : Z \to F(Z)$ is an isomorphism.
\end{theorem}

\begin{example}[Noetic Streams as Final Coalgebra]
\label{ex:noetic-streams}
For $F(D) = \Ideas \times D$, the final coalgebra is $(\Ideas^\omega, (\mathsf{head}, \mathsf{tail}))$.
\end{example}

%% ============================================================================
%% CHAPTER 5: DENOTATIONAL SEMANTICS OF TKS
%% ============================================================================

\chapter{Denotational Semantics of TKS}
\label{ch:denotational-semantics}

\begin{tcolorbox}[colback=blue!5!white,colframe=blue!75!black,title=\vnewfour]
This chapter synthesizes domain-theoretic foundations into complete
denotational semantics for TKS.
\end{tcolorbox}

\section{Semantic Domains for TKS Types}

\begin{definition}[Type Interpretation]
\label{def:type-interpretation}
The semantic domain $\sem{\tau}$ for each TKS type:
\begin{align*}
\sem{\mathsf{Idea}} &= \NoeticDomain = \Ideas_\bot \\
\sem{\mathsf{World}} &= \Worlds_\bot \\
\sem{\tau_1 \times \tau_2} &= \sem{\tau_1} \times \sem{\tau_2} \\
\sem{\tau_1 + \tau_2} &= \sem{\tau_1} \oplus \sem{\tau_2} \\
\sem{\tau_1 \to \tau_2} &= [\sem{\tau_1} \to \sem{\tau_2}] \\
\sem{\mathsf{RPM}[\tau]} &= \mathsf{Env} \to (\sem{\tau} + \mathsf{Error})_\bot
\end{align*}
\end{definition}

\begin{proposition}[Type Domains are Scott Domains]
\label{prop:type-domains-scott}
For every well-formed TKS type $\tau$, $\sem{\tau}$ is a Scott domain.
\end{proposition}

\section{Interpretation Function}

\begin{definition}[Denotational Semantics of TKS Expressions]
\label{def:denotational-semantics}
The interpretation function $\sem{-}$ is defined inductively:

\textbf{Base cases:}
\begin{align*}
\sem{x}_\rho &= \rho(x) \\
\sem{X_i}_\rho &= X_i
\end{align*}

\textbf{Noetic operators:}
\[
\sem{\noe{k}(e)}_\rho = \noe{k}_\bot(\sem{e}_\rho)
\]

\textbf{Fractals:}
\[
\sem{\langle k_1 : \cdots : k_n \rangle(e)}_\rho = (\noe{k_n}_\bot \circ \cdots \circ \noe{k_1}_\bot)(\sem{e}_\rho)
\]

\textbf{Functions:}
\begin{align*}
\sem{\lambda x.\, e}_\rho &= \lambda d.\, \sem{e}_{\rho[x \mapsto d]} \\
\sem{e_1\, e_2}_\rho &= \sem{e_1}_\rho(\sem{e_2}_\rho)
\end{align*}

\textbf{Recursion:}
\[
\sem{\mathsf{fix}\, f.\, e}_\rho = \KleeneFix(\lambda d.\, \sem{e}_{\rho[f \mapsto d]})
\]
\end{definition}

\begin{theorem}[Well-Definedness of Semantics]
\label{thm:semantics-well-defined}
For every well-typed expression $\Gamma \vdash e : \tau$:
\begin{enumerate}[label=(\roman*)]
  \item $\sem{e}_\rho$ is defined for all consistent environments.
  \item $\sem{e}_\rho \in \sem{\tau}$.
  \item The function $\rho \mapsto \sem{e}_\rho$ is Scott-continuous.
\end{enumerate}
\end{theorem}

\section{Adequacy}

\begin{theorem}[Adequacy]
\label{thm:adequacy}
For every closed expression $\vdash e : \tau$:
\[
e \bigstep v \quad \Rightarrow \quad (v, \sem{e}) \in \mathcal{R}_\tau
\]
where $\mathcal{R}_\tau$ is the adequacy relation.
\end{theorem}

\section{Full Abstraction}

\begin{theorem}[Full Abstraction for First-Order TKS]
\label{thm:full-abstraction-first-order}
The denotational semantics is fully abstract for first-order TKS:
\[
e_1 \approx_\tau e_2 \quad \Leftrightarrow \quad \sem{e_1} = \sem{e_2}
\]
\end{theorem}

\section{Connection to v7.3 Operational Semantics}

\begin{proposition}[Coherence with v7.3]
\label{prop:coherence-v73}
The denotational semantics is compatible with v7.3:
\begin{enumerate}[label=(\roman*)]
  \item Noetic operators agree with v7.0 definitions.
  \item Fractals agree with v7.0 fractal application.
  \item RPM Monad corresponds to v7.0 state transformers.
  \item Transfinite fractals correspond to domain-theoretic directed suprema.
\end{enumerate}
\end{proposition}

\begin{theorem}[Semantic Unification]
\label{thm:semantic-unification}
The domain-theoretic semantics of v7.4 provides a unified foundation for
v7.0--v7.3 semantics. Each previous version embeds faithfully into the
Scott domain framework.
\end{theorem}

\begin{tcolorbox}[colback=green!5!white,colframe=green!75!black,title=\textbf{Cross-Track Connections}]
\textbf{To Part V (Math)}: Domain equations support extended transfinite fractals.

\textbf{To Part VII (Higher Categories)}: Initial algebra/final coalgebra duality
extends to $(\infty,1)$-categories.

\textbf{To Part VIII (Geometry)}: Continuous lattices carry natural topological
structure connecting to noetic manifolds.
\end{tcolorbox}
